\section{Kubo 电导率}

Mahan 第 8 章把线性响应落到电流--电流关联，得到光学与直流电导。取 $e>0$，$\hbar=1$。

\subsection{电流算符}

均匀规范下，
\begin{equation}
  \mathbf{j}
  =\frac{e}{m L^d}
  \sum_{\mathbf{k}\sigma}
  \hbar\mathbf{k}\,
  c_{\mathbf{k}\sigma}^\dagger c_{\mathbf{k}\sigma}
  -\frac{e^2 n}{m}\mathbf{A}
  \quad
  \text{（顺磁 $+$ 抗磁）}.
\end{equation}
线性响应只保留对 $\mathbf{A}$ 的一阶。

\subsection{Kubo 公式}

诱导电流 $\langle j_\alpha\rangle=\chi_{j_\alpha j_\beta}^R(\omega)A_\beta(\omega)+\text{抗磁项}$。电导率张量
\begin{equation}
  \sigma_{\alpha\beta}(\omega)
  =\frac{\mathrm{i}e^2 n}{m\omega}\delta_{\alpha\beta}
  +\frac{\mathrm{i}}{\omega}
  \Pi_{\alpha\beta}^R(\omega),
\label{eq:kubo_sigma}
\end{equation}
其中流--流推迟关联
\begin{equation}
  \Pi_{\alpha\beta}(t)
  =-\mathrm{i}\theta(t)
  \bigl\langle
  [j_\alpha^{(p)}(t),j_\beta^{(p)}(0)]
  \bigr\rangle
\end{equation}
只用顺磁电流。规范不变性要求 $\omega\to0$ 时抗磁项与 $\Pi(0)$ 相消，绝缘体 $\sigma(0)=0$，金属留下 Drude 权重。

等价的密度响应形式（连续性）把纵向电导与 $\chi_{nn}$ 联系起来：
\begin{equation}
  \sigma_L(q,\omega)
  =\frac{\mathrm{i}e^2\omega}{q^2}\chi_{nn}(q,\omega).
\end{equation}
$q\to0$ 即光学纵向电导。

\subsection{气泡近似与 Drude}

无顶角修正时，$\Pi$ 由两个 $G$ 的气泡给出。弹性杂质、弛豫时间 $\tau$ 下
\begin{equation}
  \sigma(\omega)
  =\frac{ne^2\tau/m}{1-\mathrm{i}\omega\tau},
\label{eq:drude}
\end{equation}
直流 $\sigma_{\mathrm{dc}}=ne^2\tau/m$。电声散射使 $1/\tau(T)$ 在高温 $\propto T$，低温 $\propto T^5$（布洛赫--格林艾森，需形变势与声子占用）。

\subsection{光学定理与 \texorpdfstring{$f$}{f}-sum}

$\mathrm{Re}\,\sigma(\omega)$ 为吸收。对所有频率积分
\begin{equation}
  \int_0^{\infty}\mathrm{d}\omega\,
  \mathrm{Re}\,\sigma(\omega)
  =\frac{\pi n e^2}{2m},
\end{equation}
与密度通道 f-sum \eqref{eq:fsum} 同一守恒律。相互作用可把光谱重量从 Drude 峰转移到中红外，但积分不变。

\begin{note}
完整的“ladder + 杂质平均”给出 Boltzmann 方程的图推导（Mahan §8.1）。顶角修正在 conserving 近似中不可少，否则破坏电荷守恒。
\end{note}
